\documentclass[10pt,a4paper]{amsart}
\usepackage[margin=19mm]{geometry}
\usepackage{amsmath,amssymb,mathtools}
\usepackage[scr=rsfs,cal=euler]{mathalfa}
\usepackage{xfrac}
\usepackage[unicode,hidelinks,bookmarks=false]{hyperref}
\newcommand{\ie}{i.e.\ }
\newcommand{\cf}{cf.\ }
\newcommand{\resp}{respectively\ }
\newcommand{\Subsection}{\subsection}
\providecommand{\nobreakdash}{\nobreak\hbox{-}}
\setlength{\emergencystretch}{2em}
\multlinegap0pt
\usepackage{amssymb}
\newtheorem{lemma}{Lemma}
\title{Mean Values and Normal Order of a Successive-Iterate Ratio of Dedekind's Arithmetic Function}
\author{Aimin Guo$^{*}$}
\date{Source: arXiv:2609.06333v1 (CC BY 4.0)}
\begin{document}
\maketitle

\noindent\emph{Source and licence.} Mean Values and Normal Order of a Successive-Iterate Ratio of Dedekind's Arithmetic Function. Version arXiv:2609.06333v1 (CC BY 4.0), distributed by arXiv under the Creative Commons Attribution 4.0 International licence, http://creativecommons.org/licenses/by/4.0/, as recorded in the arXiv licence metadata for this exact version and shown on the abstract page https://arxiv.org/abs/2609.06333v1. Full licence notice: \texttt{licenses/arxiv-2609.06333-CC-BY-4.0.txt}.\par\smallskip

\noindent\textit{Editorial scope.} Counted unit: the article's lemma BT-lemma (source0.tex lines 341--350). The article's complete original proof of it is delivered with it (source0.tex lines 352--412, one \texttt{proof} environment of 61 lines). No mathematics is written, repaired or re-derived: the statement and the proof are the article's own lines, in the article's own notation, and no retained line is substituted or re-flowed.  No further source-local context is needed: every cross-reference of every retained line resolves inside the two delivered spans.  Nothing was omitted from any retained span, so the package carries no omission record.  The retained lines cite 1 works, whose entries are transcribed verbatim from the article's own bibliography source in the e-print and delivered with short editorial labels recorded in the item's reference mapping.  No retained line contains a figure environment, an image-inclusion command or drawing code, so no image is delivered and the package declares no asset dependency.  Editorial formatting only, in addition: an AMS \texttt{amsart} preamble that restates the article's own declarations and macros used by the retained lines, namely the environments \texttt{lemma} on separate counters, together with the standard packages those lines need.  The retained environments print in this document as Lemma 1 in order of appearance, whereas the article prints the same items with the numbering of its own full document.  The complete native source of the article as distributed in the arXiv e-print for this exact version is bundled unchanged as \texttt{sources/209525/source0.tex}.
\begin{lemma}\label{BT-lemma}
Let \(x\) be sufficiently large. Uniformly for primes \(q<x\),
\[
\sum_{\substack{r\le x\\r\equiv-1\pmod q}}
\frac1r
\ll
\frac{\log_2x}{q},
\]
where \(r\) runs over primes.
\end{lemma}

\begin{proof}
Let
\[
\pi(t;q,-1)
:=
\#\{r\le t:r\ \text{is prime},\ r\equiv-1\pmod q\}.
\]
If \(r\le2q\) and \(r\equiv-1\pmod q\), then
\(r=kq-1\) for some positive integer \(k\), and hence \(k=1\) or \(2\).
Thus
\[
\sum_{\substack{r\le2q\\r\equiv-1\pmod q}}\frac1r
\le
\frac1{q-1}+\frac1{2q-1}
\ll\frac1q.
\]
If \(x\le2q\), the desired estimate follows immediately. Hence we may
assume that \(x>2q\).

For \(t\ge2q\), since \((-1,q)=1\), the Brun--Titchmarsh inequality
\cite[Theorem~3.9]{MontgomeryVaughan} gives
\[
\pi(t;q,-1)
\ll
\frac{t}{\phi(q)\log(t/q)}
\ll
\frac{t}{q\log(t/q)},
\]
where we used \(\phi(q)=q-1\asymp q\).

By Abel summation,
\[
\begin{aligned}
\sum_{\substack{2q<r\le x\\r\equiv-1\pmod q}}\frac1r
&=
\frac{\pi(x;q,-1)}x
-\frac{\pi(2q;q,-1)}{2q}
+
\int_{2q}^x\frac{\pi(t;q,-1)}{t^2}\,dt\\
&\ll
\frac1q+
\frac1q\int_{2q}^x
\frac{dt}{t\log(t/q)}.
\end{aligned}
\]
With \(u=t/q\),
\[
\int_{2q}^x\frac{dt}{t\log(t/q)}
=
\int_2^{x/q}\frac{du}{u\log u}
=
\log\log(x/q)-\log\log2
\ll\log_2x.
\]
Combining this with the contribution from \(r\le2q\), we obtain
\[
\sum_{\substack{r\le x\\r\equiv-1\pmod q}}\frac1r
\ll
\frac{\log_2x}{q}.
\]
\end{proof}

\begin{thebibliography}{99}
\bibitem[Montgo]{MontgomeryVaughan} H. L. Montgomery and R. C. Vaughan, \emph{Multiplicative Number Theory I: Classical Theory}, Cambridge Studies in Advanced Mathematics, Vol.~97, Cambridge University Press, Cambridge, 2006.
\end{thebibliography}
\end{document}
